\section{部署与信任}

\subsection{评估栈与监控指标}

部署后需要把训练阶段的 evidence 连接到生产指标：
\begin{itemize}
    \item Benchmark（HumanEval、LongBench）确保能力；
    \item Guardrail（Hallucination、Forbidden topics）确保安全；
    \item Ops metrics（Rejection rate、Latency P99）确保体验。
\end{itemize}

\begin{table}[H]
\centering
\begin{tabular}{p{0.25\textwidth} p{0.55\textwidth}}
\toprule
类别 & 典型指标 \\
\midrule
Benchmark & HumanEval、TruthfulQA、Domain-specific tests \\
Safety & Red team rejection、prompt-based filter precision \\
Ops & Rejection rate、Latency P99、Human override count \\
\bottomrule
\end{tabular}
\caption{模仿学习部署后的监控栈}
\end{table}

\subsection{幻灯片与视觉证据}

\begin{figure}[H]
\centering
\includegraphics[width=0.8\textwidth]{slide-02-05.jpg}
\caption{Slide 02-05 展示了演示数据到监控的闭环（data → policy → ops board）。}
\end{figure}

\begin{figure}[H]
\centering
\includegraphics[width=0.8\textwidth]{slide-02-10.jpg}
\caption{Slide 02-10 把 evaluation stack 与 governance checkpoint 合并在一个矩阵上，强调 evidence-driven rollout。}
\end{figure}

\begin{importantbox}{可视化仪表盘的作用}
把 evidence stacking 可视化成幻灯片上的热力地图，可以让 PM、法律和 SRE 在同一个数据视角下讨论风险，而不是靠零散的 document。
\end{importantbox}

\subsection{反馈与治理}

遵循 Observe-Assess-Act-Document 模式：
\begin{itemize}
    \item Observe：记录 rollout 表现、reward drift；
    \item Assess：用 evidence matrix 检查 guardrail； 
    \item Act：human-in-the-loop 开启 mitigation；
    \item Document：把决策写入 runbook，供后续 audit。
\end{itemize}

\begin{warningbox}{记录比修复更重要}
没有 runbook 的反馈 loop，团队很难记住为什么在某个 prompt 上 rollback。Archit 建议：每次 mitigation 要附带 drift snapshot、human override log 和 lesson learned。
\end{warningbox}

\subsection{本章小结}

部署的信任来自 data、policy、ops 之间的闭环。把幻灯片中的 evidence dashboard 变成实际的监控、治理表格，才能让模仿学习系统在 production 中持续稳定。

